\section{ConVex optimization-based Steady-state Tracking Error Minimization (CV-STEM)}\label{chap:CV-STEM}

ConVex optimization-based Steady-state Tracking Error Minimization (CV-STEM) and Neural Contraction Metric (NCM) will be reviewed in this chapter and the next chapter, respectively.
H. Tsukamoto and S.-J. Chung proposed a method called ConVex optimization-based Steady-state Tracking Error Minimization (CV-STEM) in \cite{Tsukamoto:2021ae} 

First, let us consider the following system:
\begin{equation}
   \ddtt\bdx 
   =
   \bdf(\bdx(t), t)
   + 
   \bdd(t)
   ,
   \label{eq:cv:sys}
\end{equation}
where $\bdx \in \mathbb{R}^n$ is the state vector, $\bdf: \mathbb{R}^n \times \mathbb{R}^+ \to \mathbb{R}^n$ is a nonlinear function, and $\bdd(t) \in \mathbb{R}^n$ is the bounded disturbance vector such that $\norm{\bdd} \leq \overline{d}\in\R_{>0}$.
Ignoring $\bdd$, according to the Theorem \ref{thm:ctrac:main}, the condition \eqref{eq:ctrac:ctrac:metric:M} for the system \eqref{eq:cv:sys} to be contracting is given by
\begin{equation}
    \ddtt\mm{M}
    +
    \mysym(\mm{M}\tpfpx)
    \le
    -2\alpha\mm{M}
    ,
    \ 
    \forall \bdx \in \R^n, t\in\R_{>0}
    ,
    \label{eq:cv:ctrac:condition}
\end{equation}
where $\alpha \in \R_{>0}$ is the contraction rate, and $\mm{M} \in \R^{n\times n}$ is the contraction metric which satisfies $\underline{m}\mm{I} \leq \mm{M} \leq \overline{m}\mm{I}$, where $\underline{m}, \overline{m} \in \R_{>0}$ are the lower and upper bounds of the contraction metric, respectively.
Then, the smallest path integral is exponentially bounded, thereby yielding a bounded tube of states:
\begin{equation}
    \norm{\mv{\xi}_1(t)-\mv{\xi}_2(t)}
    -
    \tfrac{V_l(0)}{\sqrt{\underline{m}}}
    \exp(-\alpha t)
    \le 
    \tfrac{\overline{d}}{\alpha}
    \sqrt{
        \tfrac{\overline{m}}{\underline{m}}
    }
    (1-\exp(-\alpha t))
    \le
    \tfrac{\overline{d}}{\alpha}
    \sqrt{
        \tfrac{\overline{m}}{\underline{m}} 
    }
    \le
    \tfrac{\overline{d}}{\alpha}
    \tfrac{\overline{m}}{\underline{m}}
    ,
\end{equation}
\ie see, Theorem \ref{thm:ctrac:pert}.

\subsection{CV-STEM for deterministic systems}
\label{sec:cv:cvstem}

The CV-STEM is a method to sample contraction metrics using the convex optimization method.
The CV-STEM is introduced in \cite{Tsukamoto:2019aa,Tsukamoto:2021ad} to handle the control design of It\^{o} stochastic systems.
In \cite{Tsukamoto:2021ae}, the simpler formulation of CV-STEM  for deterministic systems is presented.

To formulate the convex optimization problem, we first reformulate the condition \eqref{eq:cv:ctrac:condition} into the following form:
\begin{align}
    J_{\rm NL}^*
    =
    \min_{
        \mathcal{X}>0,\mm{W}>0
    }
    \tfrac{\overline{d}}{\alpha}
    \mathcal{X}
    \\
    \text{s.t.}
    \eqref{eq:cv:ctrac:condition} 
    ,
    {\ \rm and \ }
    \underline{m}\mm{I} \leq \mm{M} \leq \overline{m}\mm{I}
    ,
    \label{eq:cv:cvstem:opt}
\end{align}

